% =====================================================================
%  deck.tex — presentation slides (beamer, 16:9).
%
%  STRUCTURE (standard talk arc):
%    background -> problem statement -> approach overview -> related work
%    -> METHOD (4 math slides: score graph / GP on the graph + coupling /
%    side information as kernels / full model + exact inference) ->
%    experiment setup (data, split, task, metrics, baselines) -> results
%    (main / example / ablations / deep baselines / replication /
%    findings) -> future work (Phase 2 plan, Phase 3 plan) -> summary.
%  Backup: robustness, per-channel, kernel comparison, Gaussian tail,
%  probes, Phase-2 pilot, downstream tasks.
%
%  Slide text is sparse: the speaker carries the narrative.
%  All numbers match docs/presentation_brief.md.  Graphics are \figslot
%  placeholders; set \showfigures to 1 to pull in files that exist in
%  ../thesis/figures/.
%
%  Compile:  tectonic deck.tex
% =====================================================================
\documentclass[aspectratio=169,11pt]{beamer}
\usetheme{default}
\usefonttheme{professionalfonts}
\setbeamertemplate{navigation symbols}{}
\setbeamertemplate{itemize items}{\textcolor{okblue}{\scriptsize$\blacksquare$}}
\setbeamertemplate{footline}{%
  \hfill\scriptsize\textcolor{muted}{\insertframenumber\,/\,\inserttotalframenumber}\hspace{2mm}\vspace{2mm}}

\usepackage{amsmath,amssymb,bm}
\usepackage{booktabs}
\usepackage{graphicx}
\usepackage{tikz}
\graphicspath{{../thesis/figures/}}

% --- palette (Okabe–Ito, matches every figure) -----------------------
\definecolor{okblue}{HTML}{0072B2}
\definecolor{okverm}{HTML}{D55E00}
\definecolor{okgreen}{HTML}{009E73}
\definecolor{okamber}{HTML}{E69F00}
\definecolor{muted}{HTML}{6B7280}
\definecolor{ink}{HTML}{1A1A1A}
\setbeamercolor{frametitle}{fg=ink}
\setbeamercolor{normal text}{fg=ink}
\setbeamerfont{frametitle}{series=\bfseries,size=\large}

% --- figure slots -----------------------------------------------------
% \figslot{file}{height}{what the graphic should show}
\def\showfigures{0}
\newcommand{\figslot}[3]{%
  \begin{center}
  \if1\showfigures\IfFileExists{../thesis/figures/#1}%
    {\includegraphics[height=#2]{#1}}%
    {\figslotbox{#1}{#2}{#3}}%
  \else\figslotbox{#1}{#2}{#3}\fi
  \end{center}}
\newcommand{\figslotbox}[3]{%
  \begin{tikzpicture}
    \node[draw=muted, dashed, thick, minimum height=#2, text width=0.72\textwidth,
          align=center, text=muted, inner sep=6pt]
      {\texttt{\detokenize{#1}}\\[3pt]\footnotesize #3};
  \end{tikzpicture}}

% --- notation shortcuts (match the thesis) ----------------------------
\newcommand{\R}{\mathbb{R}}
\newcommand{\N}{\mathcal{N}}
\newcommand{\Lg}{L_{\mathcal G}}
\newcommand{\diag}{\operatorname{diag}}

\title{\textbf{Expressive Performance as a Gaussian Process\\ on a Musical Score Graph}}
\author{Raynaldi Lalang\\[1mm]{\small Supervisor: Prof.\ Kazuyoshi Yoshii}}
\institute{Kyoto University}
\date{July 16, 2026}

\begin{document}

% ================================================================ 1
\begin{frame}
  \titlepage
\end{frame}

% ================================================================ 2  BACKGROUND
\begin{frame}{Expressive performance}
  The score specifies pitches and note values. Performers add everything else:
  \begin{itemize}
    \item timing: notes pushed and pulled around the beat;
    \item articulation: held longer or shorter than written;
    \item dynamics: shaped louder and softer across phrases.
  \end{itemize}
  \vspace{1mm}
  These deviations are the expression. Modelling them matters for performance
  analysis, pedagogy, and expressive rendering.
  \figslot{deck_concept.png}{2.6cm}{one bar as a piano-roll, with the played
    timing / duration / loudness deviating from the written values}
\end{frame}

% ================================================================ 3  PROBLEM
\begin{frame}{Problem: infer how the music was played, with reliable uncertainty}
  The goal: a probabilistic model of \textbf{every per-note expressive
  parameter a performer controls}: timing, articulation, dynamics, and, on
  continuously pitched instruments, intonation and vibrato; ultimately
  inferred from audio alone. Every estimate must carry an error bar that can
  be trusted note by note.
  \vspace{2mm}

  The first completed stage: piano, from aligned performance MIDI,
  \begin{equation*}
    y_i \;=\; \bigl[\;\underbrace{\tau_i}_{\text{onset timing}},\;\;
      \underbrace{\log r_i}_{\text{articulation}},\;\;
      \underbrace{v_i}_{\text{dynamics}}\;\bigr] \in \R^{3}
    \qquad\text{per written note.}
  \end{equation*}
  \begin{itemize}
    \item Not full transcription: the score is known.
    \item The output is a \textbf{per-note posterior}: an estimate and a
          calibrated error bar.
  \end{itemize}
\end{frame}

% ================================================================ 4  APPROACH OVERVIEW
\begin{frame}{Approach: one prior, three kinds of observation}
  One \textbf{Gaussian process (GP)} prior over per-note expressive fields, on
  a graph built from the score. Only the observation changes: per-note MIDI
  values (Phase 1), noisy targets derived from the $f_0$ curve (Phase 2), or
  the waveform itself (Phase 3). The Phase-0 music model supplies input
  features to all of them.
  \begin{center}
  {\footnotesize
  \begin{tabular}{@{}llll@{}}
  \toprule
  Phase & Observation & Channels & Status \\
  \midrule
  0 \; music model & MIDI corpora (MAESTRO) & --- & Done \\
  1 \; piano & performance MIDI, per note & $\tau, \log r, v$ & \textbf{Done} \\
  2 \; mono.\ pitched & $f_0$-derived targets, noisy & $+\,c$, vibrato & WIP \\
  3 \; waveform & the audio itself & all latent & To Do \\
  \bottomrule
  \end{tabular}}
  \end{center}
  \vspace{-1mm}
  \figslot{deck_programme.png}{2.4cm}{the shared prior with three swappable
    likelihood blocks (MIDI / $f_0$ / waveform)}
\end{frame}

% ================================================================ 5  RELATED WORK
\begin{frame}{Related work}
  \begin{columns}[T,onlytextwidth]
  \column{0.50\textwidth}
  \textbf{Expressive performance modelling}
  {\small
  \begin{itemize}
    \item rule systems: KTH performance rules (Friberg et al.\ 2006)
    \item linear basis-function models (Grachten \& Widmer 2012)
    \item deep rendering models: VirtuosoNet (Jeong et al.\ 2019);
          Oore et al.\ 2018
    \item review: Cancino-Chac\'on et al.\ 2018
  \end{itemize}
  All output one rendering or corpus statistics; none a calibrated
  per-note posterior.}
  \column{0.48\textwidth}
  \textbf{Gaussian processes and graphs}
  {\small
  \begin{itemize}
    \item \textbf{Gaussian Processes over Graphs} (Venkitaraman et al.\
          2018): outputs on graph nodes, Kronecker-coupled; the main
          foundation here
    \item graph Mat\'ern and heat kernels (Borovitskiy et al.\ 2021)
    \item multi-output coupling via ICM (Bonilla et al.\ 2008)
  \end{itemize}}
  \end{columns}
  \vspace{3mm}
  This work combines the two: a graph GP over score positions, with per-note
  side information, for calibrated performance inference.
\end{frame}

% ================================================================ 6  METHOD 1: SCORE GRAPH
\begin{frame}{Method (1): the score graph}
  Nodes are the written notes $(p_i, b_i, d_i)$: pitch, beat onset, duration.
  Edges connect notes that are close in score time and pitch:
  \begin{equation*}
    W_{ij} \;=\; \exp\!\Bigl( -\tfrac{(b_i - b_j)^2}{2\,\ell_b^2}
                             \;-\; \tfrac{(p_i - p_j)^2}{2\,\ell_p^2} \Bigr),
    \qquad
    \Lg \;=\; D - W, \quad D = \diag\Bigl(\textstyle\sum_j W_{ij}\Bigr).
  \end{equation*}
  \begin{itemize}
    \item $\ell_b$ (beats) and $\ell_p$ (semitones) are length scales; optional
          same-voice, same-chord, and voice-leading edges.
    \item Intuition: neighbouring notes should be played similarly. The
          Laplacian $\Lg$ turns that into a prior.
  \end{itemize}
  \figslot{deck_scoregraph.png}{2.8cm}{NEW: a few bars of real score as nodes,
    edges drawn by time / pitch / chord relation}
\end{frame}

% ================================================================ 7  METHOD 2: GP + COUPLING
\begin{frame}{Method (2): a GP on the graph, channels coupled}
  Eigendecompose the Laplacian, $\Lg = U \diag(\nu)\, U^{\!\top}$, and shape its
  spectrum with one parameter $s$:
  \begin{equation*}
    K_G(s) \;=\; U\, g(\nu; s)\, U^{\!\top},
    \qquad
    g(\nu; s) \;=\; \frac{1}{1 + s\,\nu}
    \quad\text{(graph Mat\'ern, } \alpha = 1\text{)} .
  \end{equation*}
  Smooth-on-the-graph components keep their variance; rough ones are shrunk.
  \vspace{1mm}

  Timing, articulation, and dynamics are correlated, so the three channels
  share the graph through a learned $3 \times 3$ matrix $B$:
  \begin{equation*}
    \operatorname{Cov}\bigl[\, y_{i,c},\; y_{j,c'} \bigr]
      \;=\; B_{cc'}\, \bigl[K_G(s)\bigr]_{ij}
    \qquad\Longleftrightarrow\qquad
    \operatorname{Cov}[\,y\,] \;=\; B \otimes K_G(s).
  \end{equation*}
  An observed channel at one note then informs the other channels at
  neighbouring notes.
\end{frame}

% ================================================================ 8  METHOD 3: SIDE INFORMATION
\begin{frame}{Method (3): side information as kernels}
  Two per-note feature sets: hand-built score features
  $X_{\mathrm{feat}} \in \R^{N \times 26}$ and embeddings
  $H \in \R^{N \times 512}$ from a small music language model trained on
  MAESTRO.
  \vspace{1mm}

  A linear mean with unknown weights, integrated out, becomes a kernel term:
  \begin{equation*}
    y \;=\; X\beta + f, \quad \beta \sim \N(0,\, c\,I)
    \qquad\Longrightarrow\qquad
    \operatorname{Cov}[\,y\,] \;=\; c\, X X^{\!\top} + \operatorname{Cov}[\,f\,].
  \end{equation*}
  \begin{itemize}
    \item So the regression weights are not fixed once across the corpus: they
          are \textbf{re-inferred for every piece}, jointly with everything
          else, and their uncertainty enters the posterior.
    \item This per-piece adaptation turns out to be the main source of the
          accuracy gain (results).
  \end{itemize}
\end{frame}

% ================================================================ 9  METHOD 4: FULL MODEL + INFERENCE
\begin{frame}{Method (4): full model and exact inference}
  \begin{equation*}
    K \;=\;
    \underbrace{\textcolor{okblue}{B \otimes K_G(s)}}_{\text{coupled graph GP}}
    \;+\; \underbrace{\textcolor{okamber}{\textstyle\sum_f \diag(c_f)\otimes X_f X_f^{\!\top}}}_{\text{side information}}
    \;+\; \underbrace{\textcolor{okverm}{\diag(\varsigma^2)\otimes I}}_{\text{noise}}
  \end{equation*}
  All hyperparameters $\theta = (s, B, \{c_f\}, \varsigma^2)$ from one exact
  evidence per piece; prediction is conjugate:
  \begin{equation*}
    \log p(\tilde y_o \mid \theta) = \log \N\!\bigl(\tilde y_o;\, 0,\, K_{oo}\bigr),
    \qquad
    m = C_{\cdot o} K_{oo}^{-1} \tilde y_o, \quad
    \Sigma_y = C - C_{\cdot o} K_{oo}^{-1} C_{o\cdot},
  \end{equation*}
  with $C$ the latent covariance and $o$ the observed entries.
  \begin{itemize}
    \item No variational approximation anywhere in Phase 1.
    \item The baselines are special cases: diagonal $B$, zero $c_f$, or a
          fixed cross-piece mean recover each simpler system exactly.
  \end{itemize}
  \vspace{-2mm}
  \figslot{arch_phase1_deck.png}{1.8cm}{architecture: graph branch (blue),
    feature branch (amber), likelihood (vermillion)}
\end{frame}

% ================================================================ 10  EXPERIMENT SETUP
\begin{frame}{Experiment setup}
  \begin{itemize}
    \item \textbf{Data}: ASAP piano (aligned score + performance MIDI);
          performances that appeared in pretraining are removed
          ($1036 \to 653$).
    \item \textbf{Split}: 40 train / 30 validation / 20 test
          (test evaluated once, claims fixed beforehand).
    \item \textbf{Task}: hide a fraction of notes, predict them from the
          visible rest; the masking rate is itself varied (results).
    \item \textbf{Metrics}, per channel: RMSE for accuracy; NLL and
          coverage for the error bars. Comparisons are paired per piece.
  \end{itemize}
\end{frame}

% ================================================================ 10b  SYSTEMS
\begin{frame}{Systems compared, and what each one tests}
  \begin{center}
  {\footnotesize
  \begin{tabular}{@{}lll@{}}
  \toprule
  System & What changes & Question it answers \\
  \midrule
  \textbf{proposed} & full covariance & \\
  no graph & $K_G \to I$ & does the score graph matter? \\
  no music model & embedding kernel removed & do learned features matter? \\
  fixed mean & per-piece weighting $\to$ one & does per-piece adaptation \\
   & cross-piece regression & matter? \\
  mean + smoothing & regression mean, then residuals & does joint modelling beat \\
  (plain / best) & graph-smoothed per channel & the standard two-step recipe? \\
  NN with uncertainty & GP replaced by an MLP predicting & is the GP needed at all? \\
  (single / ensemble) & mean and variance per note & \\
  independent / temporal & no edges / time-only edges & are score-graph edges needed? \\
  \bottomrule
  \end{tabular}}
  \end{center}
  \vspace{1mm}
  Every system uses the same input features and the same protocol; all but
  the NN are exact special cases of the proposed model.
\end{frame}

% ================================================================ 11a  EXPERIMENT OVERVIEW
\begin{frame}{Experiments run}
  \begin{center}
  {\footnotesize
  \begin{tabular}{@{}lll@{}}
  \toprule
  Experiment & Question & Where \\
  \midrule
  main comparison (5 systems) & does the joint model win? & next slide \\
  masking-rate sweep, 50\% to leave-one-out & does the rate matter? & results \\
  ablations (graph / embeddings / fixed mean) & where do gains come from? & results \\
  neural-network baselines & is the GP needed? & results \\
  kernel and graph variants (7 kernels, 4 graphs) & which structure matters? & results + backup \\
  music-theory features + embedding probes & does tonality help? & results \\
  seed robustness (12 mask seeds) & seed artifacts? & backup \\
  replication on 30 further unseen pieces & does it hold on new data? & backup \\
  statistical battery (5 test types, BH, clustering) & test-choice artifacts? & backup \\
  downstream tasks (6) & what is it useful for? & backup \\
  Phase-2 synthetic pilot & does the extension work? & future work \\
  \bottomrule
  \end{tabular}}
  \end{center}
\end{frame}

% ================================================================ 11  RESULTS: MAIN
\begin{frame}{Results on the test set, per channel}
  \begin{center}
  {\small RMSE / NLL per channel (coverage@90\% is 0.90--0.95 everywhere) \\[1mm]
  \begin{tabular}{@{}lccc@{}}
  \toprule
  System & timing $\tau$ & articulation $\log r$ & dynamics $v$ \\
  \midrule
  \textbf{proposed model} & \textbf{0.227} / $-0.54$ & \textbf{0.605} / $\bm{+0.83}$ & \textbf{0.076} / $\bm{-1.19}$ \\
  \quad no graph & 0.236 / $-0.52$ & 0.656 / $+0.93$ & 0.084 / $-1.09$ \\
  \quad no music model & 0.230 / $-0.51$ & 0.621 / $+0.85$ & 0.079 / $-1.13$ \\
  mean + smoothing (best) & 0.237 / $\bm{-0.72}$ & 0.632 / $+0.89$ & 0.082 / $-1.10$ \\
  mean + smoothing (plain) & 0.237 / $-0.73$ & 0.648 / $+0.92$ & 0.084 / $-1.06$ \\
  \bottomrule
  \end{tabular}}
  \end{center}
  \begin{itemize}
    \item Best RMSE in every channel; best NLL in articulation and dynamics.
    \item Timing NLL favours mean + smoothing here: traced to one piece with
          extreme timing outliers under the Gaussian likelihood (backup).
    \item Preregistered pooled claims: $\Delta$RMSE vs the best baseline
          $\bm{-0.0137^{*}}$; graph's calibration contribution
          $\Delta$NLL $\bm{-0.0736^{*}}$.
  \end{itemize}
\end{frame}

% ================================================================ 11b  MASKING RATES
\begin{frame}{Results at different masking rates (validation set)}
  \figslot{masksweep_perchannel_dev.png}{3.6cm}{per channel: RMSE vs hidden
    fraction, 50\% to 10\% plus leave-one-out, three systems (exists)}
  \begin{itemize}
    \item Articulation and dynamics: the proposed model is best at every
          rate, and the gaps grow as more context is observed.
    \item Timing: the three systems are within noise of each other at every
          rate.
    \item The leave-one-out limit agrees; conclusions do not depend on the
          40\% operating point.
  \end{itemize}
\end{frame}

% ================================================================ 12  RESULTS: EXAMPLE
\begin{frame}{Results: posterior on one piece (validation set)}
  \figslot{posterior_example_dev.png}{5.6cm}{one piece, 40\% hidden:
    posterior mean, 90\% band, observed notes, hidden truths as open circles
    (exists)}
  \vspace{-1mm}
  \centering\small Hidden true values fall inside the 90\% band at 49/50,
  49/50, 46/50 here.
\end{frame}

% ================================================================ 13  RESULTS: ABLATIONS
\begin{frame}{Results: where the gains come from (validation set)}
  Paired per-channel $\Delta$RMSE when an ingredient is removed
  (negative = the ingredient helps):
  \begin{center}
  {\small
  \begin{tabular}{@{}lccc@{}}
  \toprule
  Ingredient removed & timing $\tau$ & articulation $\log r$ & dynamics $v$ \\
  \midrule
  the graph & $-0.0015^{*}$ & $-0.0282^{*}$ & $-0.0104^{*}$ \\
  the embeddings & $-0.0006$ & $-0.0147^{*}$ & $-0.0028^{*}$ \\
  \bottomrule
  \end{tabular}}
  \end{center}
  \vspace{1mm}
  \begin{itemize}
    \item Fixing a cross-piece mean inside the GP loses almost nothing
          ($0.390$ vs $0.388$ pooled): channel coupling alone adds little.
          The accuracy gain comes from re-weighting the features per piece.
    \item The graph's calibration gain is mostly articulation (coverage
          $0.784 \to 0.91$ over a bare mean).
    \item Timing barely moves anywhere: rubato is absorbed by the tempo curve
          before $\tau$ is even defined.
  \end{itemize}
\end{frame}

% ================================================================ 14  RESULTS: DEEP BASELINES
\begin{frame}{Results: neural-network baselines (validation set)}
  An MLP that predicts a mean and a variance for every note and channel,
  trained with Gaussian likelihood on the training pieces; and an ensemble of
  five such MLPs. Same input features as the proposed model.
  \begin{center}
  {\small RMSE / NLL per channel \\[1mm]
  \begin{tabular}{@{}lccc@{}}
  \toprule
  System & timing $\tau$ & articulation $\log r$ & dynamics $v$ \\
  \midrule
  \textbf{proposed model (GP)} & \textbf{0.152} / $\bm{-0.82}$ & \textbf{0.601} / $\bm{+0.84}$ & \textbf{0.072} / $\bm{-1.23}$ \\
  MLP ensemble ($5\times$) & 0.163 / $-0.53$ & 0.757 / $+1.10$ & 0.095 / $-0.89$ \\
  MLP (single) & 0.164 / $-0.48$ & 0.762 / $+1.11$ & 0.096 / $-0.84$ \\
  \bottomrule
  \end{tabular}}
  \end{center}
  \begin{itemize}
    \item The GP is better in every channel on both metrics.
    \item The MLPs cannot adapt to the individual piece and have no graph.
    \item Caveat: the music-model features are frozen; end-to-end fine-tuning
          was not tested.
  \end{itemize}
\end{frame}

% ================================================================ 16  RESULTS: FINDINGS
\begin{frame}{Results: music-theory experiments (validation set)}
  What musical structure actually helps? Five experiments:
  \begin{enumerate}
    \item \textbf{Tonal pitch metric}: replace semitone distance in the graph
          with a pitch-class metric based on the circle of fifths:
          \emph{worse} ($+0.009^{*}$).
    \item \textbf{Chord + voice-leading edges}: add same-onset and stepwise
          edges: helps the standard construction ($-0.008^{*}$); redundant in
          the full model once the embeddings are present.
    \item \textbf{Sustain-overlap edges} (notes sounding through another's
          onset): small consistent trend, not significant.
    \item \textbf{14 music-theory features} (key and scale degree, metrical
          weight, chord dissonance, phrase boundaries, motif repetition):
          no measurable gain ($-0.002$ ns).
    \item \textbf{Probes of the embeddings} (linear and nonlinear): no tonal
          information; bass membership (AUC 0.99), meter, phrase, and
          register are encoded.
  \end{enumerate}
  \vspace{1mm}
  Tonality does not predict these channels; rhythm, voicing, and register do.
\end{frame}

% ================================================================ 17  FUTURE: PHASE 2
\begin{frame}{Future work: Phase 2, intonation and vibrato}
  Continuously pitched instruments; the channel set grows:
  \begin{equation*}
    y_i = [\,\tau_i, \log r_i, \ell_i, c_i, f_i^{\mathrm{vib}}, \gamma_i,
    \delta_i^{\mathrm{vib}}\,], \qquad
    c_i(t) = c_i + \gamma_i \sin\bigl(2\pi f_i^{\mathrm{vib}}
    (t - \delta_i^{\mathrm{vib}})\bigr) + \epsilon .
  \end{equation*}
  \textbf{Pipeline}: $f_0$ tracking $\to$ cents curve $\to$ per-note joint
  least-squares fit (with parameter variances) $\to$ the same GP, now with
  per-(note, channel) masks and heteroscedastic noise.
  \begin{itemize}
    \item Built and tested: the cell-mask posterior (exact), the estimator,
          and a synthetic pilot; vibrato parameters of notes too short to
          estimate are recovered from the graph (coverage 0.87 to 0.93).
    \item Planned data: URMP (per-instrument audio + scores); solo vocal
          corpora as secondary.
    \item Planned evaluation: same held-out imputation, calibration primary;
          a fresh preregistered test set before any claim.
    \item Main risk: audio-to-score alignment error enters $\tau$ directly and
          is correlated along time; mitigation is to model it in the noise or
          exclude $\tau$ from Phase-2 claims.
  \end{itemize}
\end{frame}

% ================================================================ 18  FUTURE: PHASE 3
\begin{frame}{Future work: Phase 3, waveform likelihood}
  Replace per-note targets with the recorded audio itself:
  \begin{equation*}
    x \;=\; \Phi(z)\, a + \epsilon, \qquad
    p(x \mid z) \;=\; \N\!\bigl(x;\; \Phi\mu_a,\;
    \Phi\Sigma_a\Phi^{\!\top} + K_n\bigr),
  \end{equation*}
  where $z$ are the per-note expressive fields (under the same graph prior),
  $\Phi(z)$ is a differentiable harmonic synthesis basis, and the amplitudes
  $a$ integrate out in closed form.
  \begin{itemize}
    \item Already implemented and tested: $\Phi$ and the amplitude
          marginalisation.
    \item Open: inference over $z$. There are no per-note observations, so
          conjugacy is lost; plan is Laplace or variational inference,
          initialised from the score.
    \item Planned evaluation: synthetic audio first (parameter recovery,
          identifiability), then monophonic recordings;
          posterior-predictive resynthesis as the end-to-end check.
    \item Open question: how to measure calibration without per-note ground
          truth.
  \end{itemize}
\end{frame}

% ================================================================ 19  SUMMARY
\begin{frame}{Summary}
  \begin{itemize}
    \item One GP prior on the score graph covers the whole program; only the
          observation model changes between phases.
    \item Phase 1: accuracy gain confirmed on a one-shot test with the
          claims fixed beforehand; the graph's calibration role confirmed
          (and stable on 30 further unseen pieces; backup).
    \item Where the gains come from is measured, not assumed: per-piece
          feature weighting for accuracy, the graph for calibration.
    \item Negative results are kept: tonality does not predict these
          deviations.
    \item Phase 2 machinery is built and validated on synthetic ground truth;
          Phase 3 is scoped with its conjugate core tested.
  \end{itemize}
\end{frame}

% ================================================================ APPENDIX
\appendix
\begin{frame}[plain]
  \centering\Large\textbf{Backup}
\end{frame}

% ---------------------------------------------------------------- A1
\begin{frame}{Backup: statistical robustness (validation set)}
  \begin{itemize}
    \item Main contrast per channel: $\tau\ -0.009^{*}$,
          $\log r\ -0.036^{*}$, $v\ -0.006^{*}$. Variance-standardized pooling
          strengthens the result (27/30 pieces).
    \item Percentile, basic, and BCa bootstrap, Wilcoxon signed-rank, and
          exact sign tests agree on every key contrast.
    \item All load-bearing starred contrasts survive Benjamini--Hochberg at
          $q = 0.05$ over 19 development contrasts.
    \item Composer-clustered bootstrap keeps every key contrast significant.
          12 mask seeds move pooled numbers by at most 0.003 RMSE.
          Piece-level coverage never falls below 0.897.
  \end{itemize}
\end{frame}

% ---------------------------------------------------------------- A2
\begin{frame}{Backup: per-channel ablation ladder (validation set)}
  \begin{center}
  {\small
  \begin{tabular}{@{}lcccccc@{}}
  \toprule
   & \multicolumn{2}{c}{Timing $\tau$} & \multicolumn{2}{c}{Artic.\ $\log r$} & \multicolumn{2}{c}{Loudness $v$} \\
  \cmidrule(lr){2-3}\cmidrule(lr){4-5}\cmidrule(lr){6-7}
  System & RMSE & cov & RMSE & cov & RMSE & cov \\
  \midrule
  coregionalized GP, zero mean & 0.156 & 0.94 & 0.676 & 0.91 & 0.080 & 0.92 \\
  + features (no graph) & 0.154 & 0.95 & 0.626 & 0.91 & 0.088 & 0.90 \\
  + features, full & 0.153 & 0.95 & 0.615 & 0.91 & 0.075 & 0.91 \\
  \textbf{+ embeddings (proposed)} & \textbf{0.152} & 0.95 & \textbf{0.601} & 0.91 & \textbf{0.072} & 0.92 \\
  \bottomrule
  \end{tabular}}
  \end{center}
  \small The calibration gain sits in articulation, the accuracy gain in
  loudness. Timing is near-static; rubato is absorbed by the tempo warp.
\end{frame}

% ---------------------------------------------------------------- A3
\begin{frame}{Backup: kernel comparison (validation set)}
  \begin{center}
  {\footnotesize
  \begin{tabular}{@{}lcccc@{}}
  \toprule
  Kernel & RMSE & NLL & cov@.9 & \\
  \midrule
  additive Laplacian (baseline) & 0.393 & $-0.322$ & 0.921 & \\
  graph Mat\'ern $\alpha \in \{1,2,3\}$ & 0.392--0.393 & $-0.318$..$-0.321$ & & tie \\
  diffusion / heat & 0.393 & $-0.312$ & 0.917 & tie \\
  independent (no edges) & $0.446^{*}$ & $-0.190^{*}$ & 0.919 & worst \\
  temporal chain & 0.392 & $-0.292^{*}$ & 0.916 & calib.\ worse \\
  tonal-distance metric (replaced) & $0.403^{*}$ & $-0.300^{*}$ & 0.921 & worse \\
  \textbf{chord + voice-leading (added)} & $\bm{0.385^{*}}$ & $\bm{-0.335^{*}}$ & 0.922 & only both-axes win \\
  \bottomrule
  \end{tabular}}
  \end{center}
  \figslot{spectral_overlay_dev.png}{2.2cm}{evidence-fitted filters
    $g(\nu;\hat s)$ per family, overlaid (exists)}
  \vspace{-2mm}
  {\footnotesize Two-stage regime. Under the final model the added edge
  families are redundant once the embeddings enter the kernel.}
\end{frame}

% ---------------------------------------------------------------- A4
\begin{frame}{Backup: the Gaussian tail and the Student-$t$ fix}
  \begin{itemize}
    \item Failure mode (one test piece and one dev replica): a very steady
          piece gives a tiny fitted timing scale; a few outlier notes then
          explode the pooled NLL while RMSE and coverage stay normal
          (dev replica: NLL $+34.7$, coverage 0.90).
    \item Student-$t$ timing likelihood (EM prototype, unit-tested): on the dev
          tail cell, NLL goes from $+34.7$ to $\bm{-0.96}$ under its own
          predictive. RMSE and coverage are unchanged.
    \item On tail-free pieces it costs a small premium.
    \item It is applied to no thesis number: the test set is spent, so the fix
          waits for its own preregistered confirmation.
  \end{itemize}
\end{frame}

% ---------------------------------------------------------------- A5
\begin{frame}{Backup: probing the embeddings}
  \begin{center}
  {\small
  \begin{tabular}{@{}lccl@{}}
  \toprule
  Theory column & $R^2$ (emb) & $R^2$ (25 feats) & verdict \\
  \midrule
  in-scale flag & $-0.06$ & $-0.00$ & not encoded (AUC 0.52) \\
  scale degree & $-0.05$ & $-0.00$ & not encoded \\
  mode (major/minor) & $-0.17$ & $-0.04$ & not encoded (AUC 0.48) \\
  repetition count & $-0.05$ & $-0.00$ & not encoded \\
  metrical weight & 0.14 & 0.43 & weakly encoded \\
  phrase salience & 0.48 & 0.57 & encoded \\
  bass flag & 0.73 & 0.86 & encoded (AUC 0.99) \\
  register & 0.90 & 1.00 & encoded \\
  \bottomrule
  \end{tabular}}
  \end{center}
  \small Cross-piece ridge probes, fit on head pieces, scored out-of-sample on
  dev pieces. RFF-2048 nonlinear probes give the same tonal verdicts.
\end{frame}

% ---------------------------------------------------------------- A6
\begin{frame}{Backup: Phase-2 synthetic pilot, per channel}
  \begin{center}
  {\footnotesize
  \begin{tabular}{@{}llll@{}}
  \toprule
  Question / system & $c$ (cents) & $\log\gamma$ & $\log f$ \\
  \midrule
  \textbf{imputation}: GP (as-given noise) & 3.51 / 0.89 & 0.094 / 0.86 & 0.026 / 0.90 \\
  \quad no-graph & 4.11 / 0.88 & 0.103 / 0.85 & 0.030 / 0.89 \\
  \quad estimator (oracle) & 1.34 / 0.93 & 0.041 / 0.90 & 0.007 / 0.93 \\
  \addlinespace
  \textbf{missing cells}: GP (as-given) & --- & 0.085 / 0.87 & 0.025 / 0.93 \\
  \quad no-graph & --- & 0.096 / 0.86 & 0.028 / 0.91 \\
  \quad estimator & \multicolumn{3}{l}{\emph{undefined; only the prior can answer}} \\
  \addlinespace
  \textbf{denoising}: GP (as-given) & 1.41 / 0.90 & 0.039 / 0.86 & 0.008 / 0.87 \\
  \quad raw estimator & 1.63 / 0.91 & 0.042 / 0.90 & 0.008 / 0.88 \\
  \bottomrule
  \end{tabular}}
  \end{center}
  \small Cells: RMSE / cov@90 vs known truth. Prior scales: $\log\gamma$ 0.35,
  $\log f$ 0.10. 926 missing vibrato cells across 20 pieces.
\end{frame}

% ---------------------------------------------------------------- A7
\begin{frame}{Backup: the 25 hand-built score features}
  {\small
  \begin{itemize}
    \item \textbf{Pitch (7)}: absolute pitch; piece-normalized pitch; local
          deviation from the 8-note window mean; local spread; melodic
          interval from the previous note; contour (sign of the interval);
          local pitch-vs-time trend.
    \item \textbf{Duration (3)}: log duration; piece-normalized; relative to
          the local window.
    \item \textbf{Meter (6)}: sine and cosine of the metrical phase at beat,
          half-bar, and bar periods.
    \item \textbf{Rhythm (3)}: inter-onset interval to the previous and next
          note; local note density.
    \item \textbf{Chords (2)}: chord size (notes sharing the onset); pitch
          rank within the chord.
    \item \textbf{Position (3)}: position in the piece; proximity to the
          start; proximity to the end.
    \item \textbf{Voice (1)}: normalized voice index.
  \end{itemize}}
  \small All computed from the score alone; a bias column is appended.
  The 14 optional music-theory columns (key, scale degree, dissonance,
  phrase, repetition) are off: they add nothing measurable (findings slide).
\end{frame}

% ---------------------------------------------------------------- A8
\begin{frame}{Backup: replication on 30 further unseen pieces}
  After the test run, the ablation was repeated on 30 additional pieces that
  no decision ever touched.
  \begin{center}
  {\small RMSE / coverage@90\% per channel \\[1mm]
  \begin{tabular}{@{}lccc@{}}
  \toprule
  System & timing $\tau$ & articulation $\log r$ & dynamics $v$ \\
  \midrule
  \textbf{proposed model} & \textbf{0.232} / 0.96 & \textbf{0.614} / 0.91 & \textbf{0.070} / 0.92 \\
  \quad no graph & 0.235 / 0.96 & 0.646 / 0.91 & 0.081 / 0.93 \\
  \quad no music model & 0.233 / 0.96 & 0.634 / 0.91 & 0.074 / 0.92 \\
  \bottomrule
  \end{tabular}}
  \end{center}
  \small Ordering and both ingredient effects replicate; the graph's
  calibration contribution is stable across validation, test, and these
  pieces ($-0.067$, $-0.074$, $-0.080$ NLL).
\end{frame}

% ---------------------------------------------------------------- A9
\begin{frame}{Backup: downstream tasks (validation set)}
  \begin{center}
  {\small
  \begin{tabular}{@{}llc@{}}
  \toprule
  Task & Verdict & Key number \\
  \midrule
  error spotting (anomaly) & \textcolor{okgreen}{\textbf{win}} & AUROC up to 0.995 \\
  denoising (known noise) & \textcolor{okgreen}{\textbf{win}} & cov 0.90 \\
  selective prediction & \textcolor{okgreen}{win} & confident-half RMSE 0.076 \\
  completion from an excerpt & \textcolor{okamber}{partial} & interpolation only \\
  era classification & \textcolor{okverm}{negative} & 0.20 vs 0.60 majority \\
  performer ID & \textcolor{okamber}{mixed} & smoothing removes identity \\
  \bottomrule
  \end{tabular}}
  \end{center}
  \begin{itemize}
    \item The graph helps note-level decisions, not piece-level summaries.
    \item Per-piece adaptation helps interpolation, fails extrapolation from
          a short excerpt.
  \end{itemize}
\end{frame}

\end{document}
